% GENERATED FILE. Edit canonical lessons, then run npm run generate:books.
% canonical-source-hash: fnv1a64:63d87e6a86d0f38f
% canonical-lessons: MR-C502-tatatut, MR-C502-kurbur, MR-C502-duravne, MR-C502-dhusphusne, MR-C502-abola-dharne

\chapter{Drifting Apart}
\label{ch:mr-drifting-apart}
\hlchaptermodality{502}

\begin{chapteropening}
\textbf{By the end of this chapter:} \emph{I can talk about a separation, bickering and falling out.}

Complete the last lesson of chapter 502: I can talk about a separation, bickering and falling out.
\end{chapteropening}

\section[\texorpdfstring{\mr{ताटातूट} (tāṭātūṭ)}{ताटातूट (tāṭātūṭ)}]{\mr{ताटातूट} (tāṭātūṭ) — a separation (of people)}

\label{lesson:MR-C502-tatatut}



\begin{quote}
Before the new one: say the Marathi for to pat, then the Marathi for to caress.
\end{quote}

\subsection*{You'll want to know: \mr{ताटातूट}}
\textbf{\mr{ताटातूट}} — \emph{tāṭātūṭ} — \textquotedblleft{}a separation (of people)\textquotedblright{}.

\begin{quote}
\mr{नोकरीमुळे} \mr{दोन} \mr{भावांची} \mr{ताटातूट} \mr{झाली}.
\end{quote}

\begin{quote}
\emph{nokrīmuḷe don bhāvāñcī tāṭātūṭ jhālī.}
\end{quote}

\begin{quote}
Work separated the two brothers.
\end{quote}

Feminine.

\mr{ताटातूट} \mr{होणे}, to be parted.

\begin{cousinweb}
\textbf{\mr{ताटातूट}} (\emph{tāṭātūṭ}) is a Marathi noun whose origin is not settled.
\end{cousinweb}

\begin{grammarlens}[title={What you've built}]
One more word for describing what you lived and felt.
\end{grammarlens}

\subsection*{Your turn}
\begin{itemize}
\raggedright
  \item \emph{Say it:} \emph{tāṭātūṭ}
  \item \emph{Say it:} \emph{tāṭātūṭ}, once more
  \item \emph{Say it:} \emph{nokrīmuḷe don bhāvāñcī tāṭātūṭ jhālī.} — Work separated the two brothers.
  \item \emph{Recall:} say the Marathi for to pat, then the Marathi for to caress, then say \emph{tāṭātūṭ} again
\end{itemize}

\begin{tcolorbox}[breakable,colback=teal!4,colframe=teal!35!black,title={Before you move on}]
What does \mr{ताटातूट} mean? (\textquotedblleft{}A separation (of people)\textquotedblright{}.) Say it once more.
\end{tcolorbox}
\section[\texorpdfstring{\mr{कुरबुर} (kurbur)}{कुरबुर (kurbur)}]{\mr{कुरबुर} (kurbur) — bickering}

\label{lesson:MR-C502-kurbur}



\begin{quote}
Before the new one: say the Marathi for to caress, then the Marathi for a separation.
\end{quote}

\subsection*{You'll want to know: \mr{कुरबुर}}
\textbf{\mr{कुरबुर}} — \emph{kurbur} — \textquotedblleft{}bickering\textquotedblright{}.

\begin{quote}
\mr{घरात} \mr{लहान} \mr{गोष्टींवरून} \mr{कुरबुर} \mr{चालते}.
\end{quote}

\begin{quote}
\emph{gharāt lahān goṣṭīṁvarūn kurbur cālte.}
\end{quote}

\begin{quote}
At home there is bickering over little things.
\end{quote}

Feminine.

Petty grumbles and squabbles: \mr{कुरबुर} \mr{करणे}.

\begin{cousinweb}
\textbf{\mr{कुरबुर}} (\emph{kurbur}) is a Marathi noun whose origin is not settled.
\end{cousinweb}

\begin{grammarlens}[title={What you've built}]
One more word for describing what you lived and felt.
\end{grammarlens}

\subsection*{Your turn}
\begin{itemize}
\raggedright
  \item \emph{Say it:} \emph{kurbur}
  \item \emph{Say it:} \emph{kurbur}, once more
  \item \emph{Say it:} \emph{gharāt lahān goṣṭīṁvarūn kurbur cālte.} — At home there is bickering over little things.
  \item \emph{Recall:} say the Marathi for to caress, then the Marathi for a separation, then say \emph{kurbur} again
\end{itemize}

\begin{tcolorbox}[breakable,colback=teal!4,colframe=teal!35!black,title={Before you move on}]
What does \mr{कुरबुर} mean? (\textquotedblleft{}Bickering\textquotedblright{}.) Say it once more.
\end{tcolorbox}
\section[\texorpdfstring{\mr{दुरावणे} (durāvṇe)}{दुरावणे (durāvṇe)}]{\mr{दुरावणे} (durāvṇe) — to drift apart}

\label{lesson:MR-C502-duravne}



\begin{quote}
Before the new one: say the Marathi for a separation, then the Marathi for bickering.
\end{quote}

\subsection*{You'll want to know: \mr{दुरावणे}}
\textbf{\mr{दुरावणे}} — \emph{durāvṇe} — \textquotedblleft{}to drift apart\textquotedblright{}.

\begin{quote}
\mr{शहर} \mr{बदलल्यावर} \mr{जुने} \mr{मित्र} \mr{दुरावले}.
\end{quote}

\begin{quote}
\emph{śahar badallyāvar june mitra durāvle.}
\end{quote}

\begin{quote}
After moving towns, old friends drifted apart.
\end{quote}

Past (he …): \mr{तो} \mr{दुरावला}, \emph{to durāvlā}.

From \mr{दूर}, far: to grow distant. \mr{दुरावा}, estrangement.

\begin{cousinweb}
\textbf{\mr{दुरावणे}} (\emph{durāvṇe}) is a Marathi verb built on \mr{दूर}, \textquotedblleft{}far\textquotedblright{}, from Sanskrit \emph{dūra}, \textquotedblleft{}distant\textquotedblright{} (CDIAL 6495), with \emph{-āvṇe}.
\end{cousinweb}

\begin{grammarlens}[title={What you've built}]
One more word for describing what you lived and felt.
\end{grammarlens}

\subsection*{Your turn}
\begin{itemize}
\raggedright
  \item \emph{Say it:} \emph{durāvṇe}
  \item \emph{Say it:} \emph{durāvṇe}, once more
  \item \emph{Say it:} \emph{śahar badallyāvar june mitra durāvle.} — After moving towns, old friends drifted apart.
  \item \emph{Recall:} say the Marathi for a separation, then the Marathi for bickering, then say \emph{durāvṇe} again
\end{itemize}

\begin{tcolorbox}[breakable,colback=teal!4,colframe=teal!35!black,title={Before you move on}]
What does \mr{दुरावणे} mean? (\textquotedblleft{}To drift apart\textquotedblright{}.) Say it once more.
\end{tcolorbox}
\section[\texorpdfstring{\mr{धुसफुसणे} (dhusphusṇe)}{धुसफुसणे (dhusphusṇe)}]{\mr{धुसफुसणे} (dhusphusṇe) — to seethe; to fume quietly}

\label{lesson:MR-C502-dhusphusne}



\begin{quote}
Before the new one: say the Marathi for bickering, then the Marathi for to drift apart.
\end{quote}

\subsection*{You'll want to know: \mr{धुसफुसणे}}
\textbf{\mr{धुसफुसणे}} — \emph{dhusphusṇe} — \textquotedblleft{}to seethe; to fume quietly\textquotedblright{}.

\begin{quote}
\mr{भावाने} \mr{ऐकलं} \mr{नाही} \mr{म्हणून} \mr{ती} \mr{धुसफुसत} \mr{होती}.
\end{quote}

\begin{quote}
\emph{bhāvāne aiklaṁ nāhī mhaṇūn tī dhusphusat hotī.}
\end{quote}

\begin{quote}
She was fuming because her brother did not listen.
\end{quote}

Past (he …): \mr{तो} \mr{धुसफुसला}, \emph{to dhusphuslā}.

Anger kept under the breath; \mr{धुसफूस}, simmering discontent.

\begin{cousinweb}
\textbf{\mr{धुसफुसणे}} (\emph{dhusphusṇe}) is a Marathi verb whose origin is not settled.
\end{cousinweb}

\begin{grammarlens}[title={What you've built}]
One more word for describing what you lived and felt.
\end{grammarlens}

\subsection*{Your turn}
\begin{itemize}
\raggedright
  \item \emph{Say it:} \emph{dhusphusṇe}
  \item \emph{Say it:} \emph{dhusphusṇe}, once more
  \item \emph{Say it:} \emph{bhāvāne aiklaṁ nāhī mhaṇūn tī dhusphusat hotī.} — She was fuming because her brother did not listen.
  \item \emph{Recall:} say the Marathi for bickering, then the Marathi for to drift apart, then say \emph{dhusphusṇe} again
\end{itemize}

\begin{tcolorbox}[breakable,colback=teal!4,colframe=teal!35!black,title={Before you move on}]
What does \mr{धुसफुसणे} mean? (\textquotedblleft{}To seethe; to fume quietly\textquotedblright{}.) Say it once more.
\end{tcolorbox}
\section[\texorpdfstring{\mr{अबोला} \mr{धरणे} (abolā dharṇe)}{अबोला धरणे (abolā dharṇe)}]{\mr{अबोला} \mr{धरणे} (abolā dharṇe) — to stop speaking to}

\label{lesson:MR-C502-abola-dharne}



\begin{quote}
Before the new one: say the Marathi for to drift apart, then the Marathi for to seethe.
\end{quote}

\subsection*{You'll want to know: \mr{अबोला} \mr{धरणे}}
\textbf{\mr{अबोला} \mr{धरणे}} — \emph{abolā dharṇe} — \textquotedblleft{}to stop speaking to\textquotedblright{}.

\begin{quote}
\mr{भांडणानंतर} \mr{तिने} \mr{भावाशी} \mr{अबोला} \mr{धरला}.
\end{quote}

\begin{quote}
\emph{bhāṇḍṇānantar tine bhāvāśī abolā dharlā.}
\end{quote}

\begin{quote}
After the quarrel she stopped speaking to her brother.
\end{quote}

Past, with -\mr{ने} (he …): \mr{त्याने} \mr{अबोला} \mr{धरला}, \emph{tyāne abolā dharlā}.

\mr{अबोला}, silence kept out of hurt; the person takes -\mr{शी}. \mr{अबोला} \mr{सोडणे}, to break the silence.

\begin{cousinweb}
\textbf{\mr{अबोला} \mr{धरणे}} (\emph{abolā dharṇe}) is a Marathi phrase built on \mr{अबोला}, \textquotedblleft{}not speaking\textquotedblright{}, the negative \emph{a-} with \emph{bol-}, the root of \mr{बोलणे}, \textquotedblleft{}to speak\textquotedblright{} (CDIAL 9321), and \mr{धरणे}, \textquotedblleft{}to hold\textquotedblright{}.
\end{cousinweb}

\begin{grammarlens}[title={What you've built}]
One more word for describing what you lived and felt.
\end{grammarlens}

\subsection*{Your turn}
\begin{itemize}
\raggedright
  \item \emph{Say it:} \emph{abolā dharṇe}
  \item \emph{Say it:} \emph{abolā dharṇe}, once more
  \item \emph{Say it:} \emph{bhāṇḍṇānantar tine bhāvāśī abolā dharlā.} — After the quarrel she stopped speaking to her brother.
  \item \emph{Recall:} say the Marathi for to drift apart, then the Marathi for to seethe, then say \emph{abolā dharṇe} again
\end{itemize}

\begin{tcolorbox}[breakable,colback=teal!4,colframe=teal!35!black,title={Before you move on}]
What does \mr{अबोला} \mr{धरणे} mean? (\textquotedblleft{}To stop speaking to\textquotedblright{}.) Say it once more.
\end{tcolorbox}
